% ═══════════════════════════════════════════════════════════════
% ML-08d: Musterlösung Repaso Final (Jahreswiederholung)
% Spanisch Klasse 7, Sequenz 8
% ═══════════════════════════════════════════════════════════════

\documentclass[11pt, a4paper]{article}

\newif\ifswmodus
\swmodusfalse

\usepackage[utf8]{inputenc}
\usepackage[T1]{fontenc}
\usepackage[ngerman]{babel}
\usepackage{helvet}
\renewcommand{\familydefault}{\sfdefault}

\usepackage[margin=2cm]{geometry}
\usepackage{enumitem}
\usepackage{tabularx}
\usepackage{booktabs}
\usepackage{multirow}

\usepackage{xcolor}
\usepackage{framed}
\usepackage{tcolorbox}
\tcbuselibrary{skins}

\usepackage{fancyhdr}
\usepackage{lastpage}
\usepackage{needspace}

% ═══════════════════════════════════════════════════════════════
% FARBEN
% ═══════════════════════════════════════════════════════════════

\definecolor{spanisch-color}{HTML}{c41e3a}
\definecolor{grau-color}{HTML}{888888}
\definecolor{hellgrau-color}{HTML}{f5f5f5}
\definecolor{orange-color}{HTML}{b35c00}
\definecolor{loesung-color}{HTML}{c62828}

\definecolor{sw-dark}{HTML}{000000}
\definecolor{sw-gray}{HTML}{666666}
\definecolor{sw-white}{HTML}{ffffff}

\ifswmodus
    \colorlet{spanisch}{sw-dark}
    \colorlet{grau}{sw-gray}
    \colorlet{hellgrau}{sw-white}
    \colorlet{orange}{sw-dark}
    \colorlet{loesung}{sw-dark}
\else
    \colorlet{spanisch}{spanisch-color}
    \colorlet{grau}{grau-color}
    \colorlet{hellgrau}{hellgrau-color}
    \colorlet{orange}{orange-color}
    \colorlet{loesung}{loesung-color}
\fi

\newcommand{\fachfarbe}{spanisch}

% ═══════════════════════════════════════════════════════════════
% VARIABLEN
% ═══════════════════════════════════════════════════════════════

\newcommand{\thema}{Repaso Final -- Jahreswiederholung}
\newcommand{\sequenz}{08}
\newcommand{\block}{d}
\newcommand{\fach}{Spanisch}
\newcommand{\klasse}{7}
\newcommand{\maxpunkte}{20}
\newcommand{\datum}{\today}

% ═══════════════════════════════════════════════════════════════
% KOPF- UND FUSSZEILE
% ═══════════════════════════════════════════════════════════════

\pagestyle{fancy}
\fancyhf{}
\renewcommand{\headrulewidth}{0.4pt}
\renewcommand{\footrulewidth}{0.4pt}

\lhead{\fach{} -- Klasse \klasse}
\chead{\textcolor{loesung}{\textbf{ML-\sequenz\block{} (MUSTERLÖSUNG)}}}
\rhead{\datum}

\lfoot{}
\cfoot{}
\rfoot{Seite \thepage\ von \pageref{LastPage}}

% ═══════════════════════════════════════════════════════════════
% BEFEHLE
% ═══════════════════════════════════════════════════════════════

\newcommand{\punktebox}[1]{\hfill\fbox{\small \_/#1}}
\newcommand{\lsg}[1]{\textcolor{loesung}{\textbf{#1}}}
\newcommand{\luecke}[1]{\rule{#1}{0.4pt}}

\newtcolorbox{merkkasten}[1][]{
    colback=hellgrau,
    colframe=\fachfarbe,
    colbacktitle=white,
    coltitle=\fachfarbe,
    boxrule=1pt,
    arc=0pt,
    left=8pt, right=8pt, top=8pt, bottom=8pt,
    fonttitle=\bfseries,
    attach boxed title to top left={yshift=-2mm, xshift=5mm},
    boxed title style={boxrule=0pt, colback=white},
    title=#1
}

\newtcolorbox{vokabelkasten}{
    colback=white,
    colframe=\fachfarbe,
    boxrule=0.5pt,
    arc=0pt,
    left=5pt, right=5pt, top=5pt, bottom=5pt
}

\newenvironment{aufgabe}[1]{%
    \needspace{5\baselineskip}%
    \nopagebreak[4]%
    \vspace{10pt}
    \noindent\textbf{\textcolor{\fachfarbe}{Aufgabe #1}}
    \nopagebreak[4]%
    \vspace{5pt}
    \nopagebreak[4]%
    \begin{list}{}{\leftmargin=0pt}
    \item[]
}{%
    \end{list}
}

\newcommand{\es}[1]{\textcolor{\fachfarbe}{\textbf{#1}}}

% ═══════════════════════════════════════════════════════════════
% DOKUMENT
% ═══════════════════════════════════════════════════════════════

\begin{document}

% ═══════════════════════════════════════════════════════════════
% KOPFBEREICH
% ═══════════════════════════════════════════════════════════════

\begin{center}
    {\Large\bfseries\textcolor{\fachfarbe}{Arbeitsblatt: \thema}}\\[0.3em]
    {\large\textcolor{loesung}{\textbf{--- MUSTERLÖSUNG ---}}}
\end{center}

\vspace{5pt}

\noindent
\begin{tabularx}{\textwidth}{|X|X|X|}
\hline
\textbf{Name:} \lsg{Musterschüler/in} &
\textbf{Klasse:} \klasse &
\textbf{Datum:} \datum \\
\hline
\end{tabularx}

\vspace{10pt}

% ═══════════════════════════════════════════════════════════════
% MERKKASTEN: PRETÉRITO INDEFINIDO
% ═══════════════════════════════════════════════════════════════

\begin{merkkasten}[Merkkasten: Pretérito indefinido (-ar)]
\textbf{Konjugation der -ar-Verben im Pretérito:}

\vspace{0.5em}

\begin{tabular}{|l|l|l|}
\hline
\textbf{Person} & \textbf{Endung} & \textbf{Beispiel: viajar} \\
\hline
yo & \lsg{-é} & viaj\lsg{é} \\
\hline
tú & \lsg{-aste} & viaj\lsg{aste} \\
\hline
él/ella/usted & \lsg{-ó} & viaj\lsg{ó} \\
\hline
nosotros/-as & \lsg{-amos} & viaj\lsg{amos} \\
\hline
vosotros/-as & \lsg{-asteis} & viaj\lsg{asteis} \\
\hline
ellos/ellas/ustedes & \lsg{-aron} & viaj\lsg{aron} \\
\hline
\end{tabular}

\vspace{0.5em}

\textbf{Signalwörter:} \es{ayer} (gestern), \es{la semana pasada} (letzte Woche), \es{el año pasado} (letztes Jahr)
\end{merkkasten}

\vspace{10pt}

% ═══════════════════════════════════════════════════════════════
% AUFGABE 1: Reisevokabular zuordnen
% ═══════════════════════════════════════════════════════════════

\begin{aufgabe}{1 -- Reisevokabular zuordnen}
Verbinde das spanische Wort mit der deutschen Bedeutung. \punktebox{4}
\end{aufgabe}

\begin{center}
\begin{tabular}{rcl}
\es{el viaje} & \lsg{---} & das Flugzeug \lsg{$\rightarrow$ el viaje = die Reise} \\[0.8em]
\es{el avión} & \lsg{---} & der Strand \lsg{$\rightarrow$ el avión = das Flugzeug} \\[0.8em]
\es{la playa} & \lsg{---} & die Reise \lsg{$\rightarrow$ la playa = der Strand} \\[0.8em]
\es{el hotel} & \lsg{---} & das Hotel \lsg{$\rightarrow$ el hotel = das Hotel} \\
\end{tabular}
\end{center}

\vspace{15pt}

% ═══════════════════════════════════════════════════════════════
% AUFGABE 2: Wettersätze bilden
% ═══════════════════════════════════════════════════════════════

\begin{aufgabe}{2 -- Wettersätze bilden}
Schreibe den passenden Wettersatz auf Spanisch. \punktebox{4}
\end{aufgabe}

\begin{vokabelkasten}
\textit{Wortliste:} Hace sol | Hace frío | Llueve | Nieva
\end{vokabelkasten}

\vspace{0.5em}

\noindent
a) Es regnet. $\rightarrow$ \lsg{Llueve.} \\[0.8em]
b) Es schneit. $\rightarrow$ \lsg{Nieva.} \\[0.8em]
c) Die Sonne scheint. $\rightarrow$ \lsg{Hace sol.} \\[0.8em]
d) Es ist kalt. $\rightarrow$ \lsg{Hace frío.}

\vspace{15pt}

% ═══════════════════════════════════════════════════════════════
% AUFGABE 3: Pretérito konjugieren
% ═══════════════════════════════════════════════════════════════

\begin{aufgabe}{3 -- Pretérito: Verben konjugieren}
Setze die richtige Form des Pretérito indefinido ein. \punktebox{6}
\end{aufgabe}

\noindent
a) Yo \lsg{viajé} (viajar) a España el año pasado. \\[0.8em]
b) Tú \lsg{visitaste} (visitar) la playa ayer. \\[0.8em]
c) María \lsg{nadó} (nadar) en el mar. \\[0.8em]
d) Nosotros \lsg{bailamos} (bailar) en la fiesta. \\[0.8em]
e) Vosotros \lsg{comprasteis} (comprar) recuerdos. \\[0.8em]
f) Ellos \lsg{tomaron} (tomar) el sol en la playa.

\vspace{15pt}

% ═══════════════════════════════════════════════════════════════
% AUFGABE 4: Urlaubsbericht schreiben
% ═══════════════════════════════════════════════════════════════

\begin{aufgabe}{4 -- Urlaubsbericht schreiben}
Schreibe einen kurzen Bericht über einen Urlaub (mind. 4 Sätze im Pretérito). \punktebox{6}
\end{aufgabe}

\textbf{Beispiellösung:}

\lsg{El año pasado viajé a España con mi familia.}

\lsg{Visitamos Barcelona y la playa.}

\lsg{Hizo mucho sol y nadé en el mar.}

\lsg{Compramos recuerdos y fue muy bonito.}

\vspace{0.5em}
\textit{\textcolor{grau}{Bewertung: 4 korrekte Sätze im Pretérito (4P), Wortschatz passend (1P), sprachliche Korrektheit (1P)}}

% ═══════════════════════════════════════════════════════════════
% PUNKTEÜBERSICHT
% ═══════════════════════════════════════════════════════════════

\vspace{20pt}
\hrule
\vspace{10pt}

\noindent
\begin{tabularx}{\textwidth}{|X|c|c|c|c||c|}
\hline
\textbf{Aufgabe} & 1 & 2 & 3 & 4 & \textbf{Gesamt} \\
\hline
\textbf{max. Punkte} & 4 & 4 & 6 & 6 & \textbf{20} \\
\hline
\textbf{erreicht} & \lsg{4} & \lsg{4} & \lsg{6} & \lsg{6} & \lsg{20} \\
\hline
\end{tabularx}

\end{document}
